
\subsubsection{Overview}

The methodology tests whether Mamba-2 duality equations can initialize SSM parameters from BERT attention weights such that:

\begin{enumerate}
\item \textbf{Stage 1 (Existence):} The derived parameters produce stable, non-divergent SSM outputs.
\item \textbf{Stage 2 (Mechanism):} The derived parameters reconstruct attention outputs better than random initialization.
\end{enumerate}

\subsubsection{Duality-Based Parameter Extraction}

Closed-form SSM parameter derivation from BERT attention weights is implemented following the Mamba-2 SSD framework. For each attention layer with query, key, and value projections ($W_Q, W_K, W_V \in \mathbb{R}^{768\times 768})$:

\textbf{A Matrix (State Dynamics):} $QK^T = W_Q \cdot W_K^T$ is computed and its principal components are extracted via SVD. Negative absolute values ensure stability (negative eigenvalues prevent divergence).

\textbf{B and C Matrices:} B and C are derived from the SVD components, normalized for numerical stability.

\textbf{D and $\Delta$:} Skip connection D = 0.1 and discretization step $\Delta$ = 1/$\sqrt768$.

\subsubsection{Stability Validation (h-e1)}

Parameter stability is validated by running SSM forward passes on calibration data, checking for NaN/Inf and measuring magnitude ratio relative to Transformer outputs. The threshold for acceptable magnitude ratio is set at $10\times$.

\subsubsection{Reconstruction Error Comparison (h-m1)}

Duality initialization is compared against random initialization using output Frobenius norm:

$\|SSM$ output - Attention $output\|_F$

For random initialization, A is sampled from $\mathcal{N}(0, 1/\sqrt{d_{\mathrm{state}}})$, B and C use Xavier uniform initialization, D is set to zeros, and $\Delta$ = 0.1. The same random seed (42) ensures reproducibility.
